\appendix{The prompt}

Write a piece of creative writing structured in two parts:
a short Introduction (the prelude) followed by a Story
(the fugue) that will be printed with a page-break between.
The entire piece (prelude and fugue) should feature a particular
mood, emotion or trait analogous to the musical concept of
\textit{key}.  (There are 24 keys in classical music.) The two
parts, prelude and fugue, will be linked by the key you have
chosen and need not be linked in style.

\begin{itemize}
\item Decide on a small number (from three to six) of words describing
  a mood, emotion or character trait you wish to address.  One
  suggestion is that you might choose these from listening to some
  music and recording your emotional responses to it but you can come
  up with the mood in any way you like. From now on your work will
  focus on these words and not the music if you used any.
\item Decide on a place, object, animal or person which represents
  this mood, character, emotion or traits.  This is your Subject.  (It
  is best if the subject is not musical or a piece of music but
  something quite different, though your writing may involve music if
  you wish.)
\item Write a short introduction, introducing your subject and setting
  the scene and/or providing a short back-story (the prelude). The
  prelude sets the scene but is not the main story.
\item Write a longer story about the subject (the fugue). 
\end{itemize}

Take `story' in the most general possible sense.  It
may help you (but it is not essential) to know that `Fugue' literally
means `Flight', as in `running away.'

Your work may be in any medium: prose, poetry, play, dialogue,
monologue or journalism.  It may be fact or fiction.

\appendix{The 24 musical keys}

I list here all the 24 classical \textit{keys}.  Each key has a
\textit{tonic note} (there are 12 of these) and is either
\textit{major} or \textit{minor}.  Each key has an associated number
of \textit{sharps} or \textit{flats} which indicates how it is
written.  Some keys can be written `equivalently' in either sharps or
flats, but note that the equivalence of keys written differently, for
example F\sh\ major and G\fl\ major, only means these would in
principle sound exactly the same if played on a piano keyboard or
computer. Most musicians, however, respond differently to these keys
when playing\=especially if the performance is on a different
instrument where other nuances of intonation are possible\=and so in
practice the keys could still sound different and convey a different `mood'.

In general, the more sharps or flats a key has the more unusual and
esoteric it sounds. Keys with few sharps or flats sound more `normal.'

I also include a few
words about each key that I have seen written on the internet.  (I 
don't agree with them all.)

\hfill SPK\break

\clearpage

\noindent\textbf{C major.}  No sharps or flats.
Pure, innocence, simplicity, naïvety, children's talk.  Often the
first (easiest) key we learn to play in.

\noindent\textbf{A minor.}  No sharps or flats.
Pious, feminine, tender.

\noindent\textbf{G major.}  One sharp.
Rustic, idyllic, lyrical, calm, gentle, peaceful.

\noindent\textbf{E minor.}  One sharp.
Na\"ive, innocent, lament.

\noindent\textbf{D major.}  Two sharps.
Triumphant. Confident, victorious, rejoicing, bright.

\noindent\textbf{B minor.}  Two sharps.
Quiet acceptance of fate. Submission.

\noindent\textbf{A major.}  Three sharps.
Innocent love, satisfaction, cheerfulness and trust.

\noindent\textbf{F\sh\ minor.}  Three sharps.
Gloomy. Resentment and discontent.

\noindent\textbf{E major.}  Four sharps.
Fresh, brash, new, happy, joy, laughing, pleasure.

\noindent\textbf{C\sh\ minor.}  Four sharps.
Penitential, lamentation, intimate, sighing.

\noindent\textbf{B major.}  Five sharps.
Colourful, passionate. Can incorporate anger, rage, jealousy, fury, despair.

\noindent\textbf{G\sh\ minor aka A\fl\ minor.}  Five sharps or seven flats.
Grumbling, suffocating, wailing, struggle, difficulty.

\noindent\textbf{F\sh\  major aka G\fl\ major.}  Six sharps or six flats.
Triumph over difficulty.

\noindent\textbf{D\sh\ minor aka E\fl\ minor.}  Six sharps or six flats.
Anxiety, distress,  despair,  depression, gloom.

\noindent\textbf{D\fl\ major aka C\sh\ major.}  Five flats or seven sharps.
Majestic.

\noindent\textbf{B\fl\ minor aka A\sh\ minor.}  Five flats or seven sharps.
Quaint, dressed in a garment of night. Surly and mocking.

\noindent\textbf{A\fl\ major.}  Four flats.
Pomp.  Grand. Possibly: death, grave, judgement, eternity.

\noindent\textbf{F minor.}  Four flats.
Melancholy. Sombre, emotional.

\noindent\textbf{E\fl\ major.}  Three flats.
Considered by many composers as `heroic.' Also devotion.

\noindent\textbf{C minor.}  Three flats.
Tragic, fate, enigmatic, brooding,
emotionally charged, melancholy, solemnity.

\noindent\textbf{B\fl\ major.}  Two flats.
Noble. Cheerful, clear conscience and hope.

\noindent\textbf{G minor.}  Two flats.
Sweetness, tenderness, serious, magnificent, discontent, unease, intensity.

\noindent\textbf{F major.}  One flat. 
Pastoral.  Countryside. Calm.

\noindent\textbf{D minor.}  One flat.
Melancholy, feminine, brooding humours.
